\section{Repository Observation and Field-Study Protocol}
\label{sec:field}

The simulation validates constructs. A bounded public-repository observation can
establish that named deficit classes occur outside the generator, but only
reconstruction drills can establish real-world cost magnitudes and economic
significance. We therefore report one executed observation, then specify the full
protocol---still not executed---at implementable detail. Design follows standard
guidance for case study and measurement research in software
engineering~\cite{runeson2009,wohlin2012}.

\subsection{Age-Stratified Repository Observation (Executed)}
We executed a read-only observation of the public
\href{https://github.com/argoproj/argo-cd}{\texttt{argoproj/argo-cd}} repository,
frozen at 2026-08-07 15:30 UTC. Selection uses \texttt{mergedAt}, not mutable
\texttt{updatedAt}. Four disjoint cohorts contain 50 PRs each: the latest merges
before the cutoff and the merges nearest target ages of one, three, and five years.
The collector retrieves complete, declared date windows before sorting, so search
result order does not choose the sample. We also take the 30 most recent published,
non-draft releases. This is a cross-sectional age design: cohort differences may
reflect changing contribution mix or project practice and cannot establish that a
link disappeared over time.

Under the narrow repository schema $\Sch_{\mathrm{repo}}$, external issue linkage
is (i)~a GitHub \texttt{closingIssuesReferences} relation or (ii)~a title/body token
of the form \texttt{\#123}, \texttt{owner/repo\#123}, or a full GitHub
\texttt{/issues/123} URL that resolves through the API to an Issue. Closed issues
count; PR objects, unresolved/deleted targets, bare numbers, Jira keys, PR URLs,
template placeholders, and references found only in comments do not. An
\texttt{APPROVED} review object is only a \emph{review-record presence control},
not proof of organizational authorization. Release-lineage presence is an explicit
PR or compare reference in release notes. The executable grammar and its unit tests
are released with the artifact.

% Generated by empirical/collect_argocd_repository.py
\newcommand{\ArgoPRTotal}{200}
\newcommand{\ArgoClosingPresent}{56}
\newcommand{\ArgoClosingPresentPct}{28.0\%}
\newcommand{\ArgoExplicitOnly}{39}
\newcommand{\ArgoExplicitOnlyPct}{19.5\%}
\newcommand{\ArgoNoExternal}{105}
\newcommand{\ArgoNoExternalPct}{52.5\%}
\newcommand{\ArgoApproved}{197}
\newcommand{\ArgoApprovedPct}{98.5\%}
\newcommand{\ArgoReleaseTotal}{30}
\newcommand{\ArgoReleaseLineage}{30}
\newcommand{\ArgoReleaseLineagePct}{100.0\%}
\newcommand{\ArgoRecentTotal}{50}
\newcommand{\ArgoRecentNoExternal}{30}
\newcommand{\ArgoRecentNoExternalPct}{60.0\%}
\newcommand{\ArgoOneYearTotal}{50}
\newcommand{\ArgoOneYearNoExternal}{28}
\newcommand{\ArgoOneYearNoExternalPct}{56.0\%}
\newcommand{\ArgoThreeYearTotal}{50}
\newcommand{\ArgoThreeYearNoExternal}{26}
\newcommand{\ArgoThreeYearNoExternalPct}{52.0\%}
\newcommand{\ArgoFiveYearTotal}{50}
\newcommand{\ArgoFiveYearNoExternal}{21}
\newcommand{\ArgoFiveYearNoExternalPct}{42.0\%}
\newcommand{\ArgoHumanTotal}{125}
\newcommand{\ArgoHumanNoExternal}{65}
\newcommand{\ArgoHumanNoExternalPct}{52.0\%}
\newcommand{\ArgoBotTotal}{75}
\newcommand{\ArgoBotNoExternal}{40}
\newcommand{\ArgoBotNoExternalPct}{53.3\%}
\newcommand{\ArgoFeatureBugTotal}{79}
\newcommand{\ArgoFeatureBugNoExternal}{29}
\newcommand{\ArgoFeatureBugNoExternalPct}{36.7\%}
\newcommand{\ArgoMaintenanceTotal}{88}
\newcommand{\ArgoMaintenanceNoExternal}{48}
\newcommand{\ArgoMaintenanceNoExternalPct}{54.5\%}
\newcommand{\ArgoDocumentationTotal}{33}
\newcommand{\ArgoDocumentationNoExternal}{28}
\newcommand{\ArgoDocumentationNoExternalPct}{84.8\%}

\begin{table}[t]
\centering
\caption{Public Argo CD observations at the frozen cutoff}
\label{tab:argocd-census}
\begin{tabular}{@{}p{0.58\columnwidth}rr@{}}
\toprule
Evidence observed & Count & Share \\
\midrule
Closing-issue link present & \ArgoClosingPresent/\ArgoPRTotal & \ArgoClosingPresentPct \\
Explicit non-closing issue reference only & \ArgoExplicitOnly/\ArgoPRTotal & \ArgoExplicitOnlyPct \\
No external issue reference & \ArgoNoExternal/\ArgoPRTotal & \ArgoNoExternalPct \\
$\geq 1$ approved-review record & \ArgoApproved/\ArgoPRTotal & \ArgoApprovedPct \\
Release with PR/compare reference & \ArgoReleaseLineage/\ArgoReleaseTotal & \ArgoReleaseLineagePct \\
\bottomrule
\end{tabular}
\end{table}

The first three rows of \cref{tab:argocd-census} are mutually exclusive:
\ArgoClosingPresent\ PRs retain a closing link, \ArgoExplicitOnly\ retain only a
resolved non-closing issue reference, and \ArgoNoExternal\ (\ArgoNoExternalPct)
retain neither under $\Sch_{\mathrm{repo}}$. The presence controls are high but not
artificially perfect: \ArgoApproved/\ArgoPRTotal\ PRs have an approved-review record,
and \ArgoReleaseLineage/\ArgoReleaseTotal\ releases have a change reference. Thus
the collector does not simply label every public evidence channel incomplete.

\begin{table}[t]
\centering
\caption{Missing external issue reference by descriptive stratum}
\label{tab:argocd-strata}
\begin{tabular}{@{}lrr@{}}
\toprule
Stratum & Missing/total & Share \\
\midrule
Recent & \ArgoRecentNoExternal/\ArgoRecentTotal & \ArgoRecentNoExternalPct \\
$\approx 1$ year & \ArgoOneYearNoExternal/\ArgoOneYearTotal & \ArgoOneYearNoExternalPct \\
$\approx 3$ years & \ArgoThreeYearNoExternal/\ArgoThreeYearTotal & \ArgoThreeYearNoExternalPct \\
$\approx 5$ years & \ArgoFiveYearNoExternal/\ArgoFiveYearTotal & \ArgoFiveYearNoExternalPct \\
Human-authored & \ArgoHumanNoExternal/\ArgoHumanTotal & \ArgoHumanNoExternalPct \\
Bot-authored & \ArgoBotNoExternal/\ArgoBotTotal & \ArgoBotNoExternalPct \\
Feature/bug & \ArgoFeatureBugNoExternal/\ArgoFeatureBugTotal & \ArgoFeatureBugNoExternalPct \\
Dependency/maintenance & \ArgoMaintenanceNoExternal/\ArgoMaintenanceTotal & \ArgoMaintenanceNoExternalPct \\
Documentation & \ArgoDocumentationNoExternal/\ArgoDocumentationTotal & \ArgoDocumentationNoExternalPct \\
\bottomrule
\end{tabular}
\end{table}

\Cref{tab:argocd-strata} exposes composition rather than hiding it. Bot- and
human-authored PRs have similar missing-link shares, whereas feature/bug PRs carry
more linkage than maintenance or documentation PRs. The age-cohort gradient runs
opposite a simple decay story: older cohorts retain more external issue linkage.
Because project norms and PR mix differ across eras, this is a descriptive warning
against interpreting the observation as link decay, not evidence against the
model's within-record survival mechanism.

%% Generated by empirical/validate_manual_audit.py
\newcommand{\ArgoAuditN}{50}
\newcommand{\ArgoAuditTP}{29}
\newcommand{\ArgoAuditTN}{21}
\newcommand{\ArgoAuditFP}{0}
\newcommand{\ArgoAuditFN}{0}
\newcommand{\ArgoAuditPrecision}{1.000}
\newcommand{\ArgoAuditRecall}{1.000}
\newcommand{\ArgoAuditAgreement}{1.000}

A single-reviewer, unblinded audit of \ArgoAuditN\ PRs was sampled within cohorts
using seed 20260807. Against manual title/body and resolved-target coding, the
collector produced TP~=~\ArgoAuditTP, TN~=~\ArgoAuditTN, FP~=~\ArgoAuditFP, and
FN~=~\ArgoAuditFN\ (precision~=~\ArgoAuditPrecision; recall~=~\ArgoAuditRecall).
This validates implementation of the declared linkage grammar on the audit sample;
it does not validate issue linkage as a measure of intent quality.

The observation has strict limits. A self-contained PR may document intent
adequately; an issue may exist in a private or non-GitHub tracker; a dependency
changelog may contribute a resolved link without documenting Argo CD intent;
approval may be encoded outside GitHub; and present API state cannot prove that a
link was once present and later lost. The sample is stratified and purposive, not a
prevalence sample. It establishes the occurrence of missing external issue linkage
under $\Sch_{\mathrm{repo}}$; it does not establish inadequate intent
documentation, temporal decay, organizational negligence, evidence-debt magnitude,
reconstruction cost, or financial exposure. Frozen data, classifications, audit
coding, and generated counts are released in \texttt{empirical/}.

\subsection{Instrument 1: Deficit Census}
A static analyzer over an organization's stores (forge, ITSM, CI, registry, CD,
cloud audit trail) that, for a sampled change population and a declared schema
$\Sch$, computes EDR by class and interrogative (\cref{met:edr}), the identifier
continuity rate across store boundaries (P5), and the estate's decay schedule
(retention settings, scheduled migrations, service-level bus factors from
repository history~\cite{avelino2016}). Output: the estate's deficit and decay
inventory---inputs (i) and (iii) of \cref{prop:reduction}. Non-invasive; requires
read access only.

\subsection{Instrument 2: Reconstruction Drills}
Timed, protocolized exercises measuring ERC directly (\cref{met:erc}): trained
participants answer standardized historical queries (``for this running workload,
establish what changed in window $w$, by whom, under which authorization, and
why'') using only the organization's actual stores. Recorded: person-time by
subtask, sources consulted, confidence of each recovered fact, unresolved facts,
and---where an independent witness exists (e.g., the original author)---error.
Drills are sampled across evidence ages to test IH1--IH2 (\cref{sec:interest-hyps})
and across services with differing census scores, giving the correlation between
EDR and measured ERC that a validated index (\cref{met:edi}) requires.
Ethics: participants are studied with consent under institutional review;
reconstruction targets are the organization's own records.

\subsection{Instrument 3: Longitudinal Decay Observation}
A yearly census re-run plus drill repetition on a fixed query panel. This observes
interest directly (same query, later time, cost delta), realized irrecoverability
(panel queries that become unanswerable), and the natural-experiment effects of
migrations and departures as they occur. Three years covers the retention cliffs
and a meaningful turnover fraction at typical rates~\cite{rigby2016}.

\subsection{Opportunistic Measurement}
Real incidents and audits are unplanned drills. The protocol includes a lightweight
retrospective instrument---time-to-establish-what-changed, sources used, facts
never established---attachable to existing incident-review practice~\cite{beyer2016}
so that realized liabilities (\cref{def:liability}) are captured at low cost. DORA-%
style delivery telemetry~\cite{forsgren2018,dora2019} provides covariates
(change volume, deployment frequency) for normalizing debt inflow
(\cref{met:growth}).

\subsection{What Would Refute the Framework}
The framework fails empirically if: census EDR does not predict drill ERC beyond
chance (the syntactic layer is disconnected from the economic one); measured cost
curves are smooth in evidence age with no alignment to retention/migration
schedules (against IH1 and \cref{prop:steps}); or false-fact rates in drills are
negligible even under heavy deficit (against \cref{fnd:false}'s transfer). Each
refutation would be informative, and the instruments are designed to detect them.
